\documentclass[12pt,a4paper]{article}
\usepackage[utf8]{inputenc}
\usepackage{amsmath,amssymb}
\usepackage{geometry}
\usepackage{enumitem}
\usepackage{fancyhdr}
\usepackage{lastpage}

\geometry{a4paper,top=25mm,bottom=25mm,left=25mm,right=25mm}
\setlength{\parindent}{0pt}
\setlength{\parskip}{6pt}

\pagestyle{fancy}
\fancyhf{}
\renewcommand{\headrulewidth}{0pt}
\fancyfoot[L]{Student number:\hspace{4cm}}
\fancyfoot[R]{\thepage\ of \pageref{LastPage}}

\newcommand{\N}{\mathcal{N}}
\newcommand{\KL}{\mathrm{KL}}
\newcommand{\E}{\mathbb{E}}
\newcommand{\ELBO}{\mathrm{ELBO}}

\begin{document}

\begin{center}
{\large\bfseries Technical University of Denmark}\\[4pt]
{\large Advanced machine learning (02460)}\\[4pt]
{\large Practice Exam Set A}\\[4pt]
Exam duration: 2 hours.\quad Aid: None.\\[2pt]
\textbf{IMPORTANT:} Questions are self-contained; all necessary formulas are provided.\\
Only information written on pages 2--11 will be evaluated.\\
Write your student number on all pages at the bottom-left.
\end{center}

\bigskip
\textbf{Contents}
\begin{itemize}[noitemsep]
  \item[\textbf{I}] \textbf{Generative models} \hfill 2
  \begin{itemize}[noitemsep]
    \item Exercise A: Variational autoencoder (VAE) \hfill 2
    \item Exercise B: Normalizing flows \hfill 3
    \item Exercise C: Denoising Diffusion Probabilistic Models \hfill 4
  \end{itemize}
  \item[\textbf{II}] \textbf{Representation learning} \hfill 5
  \begin{itemize}[noitemsep]
    \item Exercise D: Pull-back metrics \hfill 5
    \item Exercise E: Curve energy \hfill 6
    \item Exercise F: Riemannian integration \hfill 7
    \item Exercise G: ELBO and evidence \hfill 8
  \end{itemize}
  \item[\textbf{III}] \textbf{Graph models} \hfill 9
  \begin{itemize}[noitemsep]
    \item Exercise H: Graph Laplacians \hfill 9
    \item Exercise I: Weisfeiler-Lehman test \hfill 10
    \item Exercise J: Shallow embeddings \hfill 10
    \item Exercise K: Message passing \hfill 11
  \end{itemize}
\end{itemize}

\newpage
\section*{Part I \quad Generative models}

\subsection*{Exercise A --- Variational autoencoder (VAE)}

Consider a deep latent variable model with latent variable $z\in\mathbb{R}$ and observed variable
$\mathbf{x}\in\mathbb{R}^2$.  Let $z$ follow a standard Gaussian prior
\[
  p(z) = \N(z\mid 0, 1),
\]
and let the likelihood be
\[
  p(\mathbf{x}\mid z) = \N\!\bigl(\mathbf{x}\mid [z,\,-z]^\top,\, \mathbf{I}_2\bigr),
\]
where $\mathbf{I}_2\in\mathbb{R}^{2\times 2}$ is the identity matrix.  Consider the variational
posterior
\[
  q(z\mid\mathbf{x}) = \N\!\bigl(z\mid x_1 + x_2,\; 1\bigr).
\]
Recall that the ELBO is
\[
  \ELBO(\mathbf{x}) = \E_{z\sim q(z\mid\mathbf{x})}\!\bigl[\ln p(\mathbf{x}\mid z)\bigr]
    - \KL\!\bigl[q(z\mid\mathbf{x})\;\|\; p(z)\bigr], \tag{1}
\]
that the bivariate Gaussian density with identity covariance is
\[
  \N(\mathbf{x}\mid\boldsymbol{\mu},\mathbf{I}_2)
    = \tfrac{1}{2\pi}\exp\!\bigl(-\tfrac{1}{2}\|\mathbf{x}-\boldsymbol{\mu}\|^2\bigr), \tag{2}
\]
and that the KL between two univariate Gaussians is
\[
  \KL\!\bigl[\N(\mu_1,\sigma_1^2)\;\|\;\N(\mu_2,\sigma_2^2)\bigr]
    = \ln\frac{\sigma_2}{\sigma_1}
      + \frac{\sigma_1^2 + (\mu_1-\mu_2)^2}{2\sigma_2^2} - \frac{1}{2}. \tag{3}
\]

\textbf{Question A.1:} Assuming a one-sample Monte Carlo estimate of the reconstruction term
with sampled value $z'=1$, what is $\ELBO(\mathbf{x}')$ evaluated at $\mathbf{x}'=[2,-1]^\top$?

\begin{itemize}[noitemsep,label={}]
  \item[\textbf{A}] $\square$ \quad $\ELBO(\mathbf{x}') = -1 - \ln(2\pi)$
  \item[\textbf{B}] $\square$ \quad $\ELBO(\mathbf{x}') = -\tfrac{1}{2} - \ln(2\pi)$
  \item[\textbf{C}] $\square$ \quad $\ELBO(\mathbf{x}') = -\tfrac{3}{2} - \ln(2\pi)$
  \item[\textbf{D}] $\square$ \quad $\ELBO(\mathbf{x}') = \tfrac{1}{2} - \ln(2\pi)$
\end{itemize}

Give a brief argument for your answer.

\vspace{6cm}

\newpage
\subsection*{Exercise B --- Normalizing flows}

Consider a normalizing flow with random variables $\mathbf{u},\mathbf{x}\in\mathbb{R}^2$, where the
base distribution is $p_{\mathbf{u}}(\mathbf{u})=\N(\mathbf{u}\mid\mathbf{0},\mathbf{I}_2)$, and
the density of $\mathbf{x}$ is
\[
  p_{\mathbf{x}}(\mathbf{x}) = p_{\mathbf{u}}(\mathbf{u})\,|\det J_T(\mathbf{u})|^{-1},
  \qquad \mathbf{u} = T^{-1}(\mathbf{x}). \tag{4}
\]
Let $T = T_2\circ T_1$ with the following known values:
\begin{align*}
  T_1([0,0]^\top)=[1,0]^\top,\quad &\det J_{T_1}([0,0]^\top)=2,\\
  T_1([2,0]^\top)=[3,1]^\top,\quad &\det J_{T_1}([2,0]^\top)=4,\\
  T_2([1,0]^\top)=[1,1]^\top,\quad &\det J_{T_2}([1,0]^\top)=3,\\
  T_2([3,1]^\top)=[4,2]^\top,\quad &\det J_{T_2}([3,1]^\top)=5.
\end{align*}
Recall that $(T_2\circ T_1)^{-1}=T_1^{-1}\circ T_2^{-1}$ and
\[
  \det J_{T_2\circ T_1}(\mathbf{u}) = \det J_{T_2}(T_1(\mathbf{u}))\cdot\det J_{T_1}(\mathbf{u}). \tag{5}
\]

\textbf{Question B.1:} What is $p_{\mathbf{x}}([1,1]^\top)$?

\begin{itemize}[noitemsep,label={}]
  \item[\textbf{A}] $\square$ \quad $p_{\mathbf{x}}([1,1]^\top) = \dfrac{1}{6\pi}$
  \item[\textbf{B}] $\square$ \quad $p_{\mathbf{x}}([1,1]^\top) = \dfrac{1}{12\pi}$
  \item[\textbf{C}] $\square$ \quad $p_{\mathbf{x}}([1,1]^\top) = \dfrac{e^{-2}}{6\pi}$
  \item[\textbf{D}] $\square$ \quad $p_{\mathbf{x}}([1,1]^\top) = \dfrac{e^{-2}}{12\pi}$
\end{itemize}

Give a brief argument for your answer.

\vspace{8cm}

\newpage
\subsection*{Exercise C --- Denoising Diffusion Probabilistic Models}

Consider a DDPM with $T=5$ latent variables, where $\mathbf{x}_0\in\mathbb{R}^2$ is the data and
$\mathbf{x}_1,\ldots,\mathbf{x}_5\in\mathbb{R}^2$ are the latent variables.  Let the variance
schedule be
\[
  \beta_1 = \tfrac{1}{2},\quad \beta_2=\tfrac{1}{2},\quad
  \beta_3=\tfrac{1}{4},\quad \beta_4=\tfrac{1}{4},\quad \beta_5=\tfrac{1}{4}.
\]
Recall that the forward process satisfies
\[
  q(\mathbf{x}_{1:T}\mid\mathbf{x}_0)=\prod_{t=1}^T q(\mathbf{x}_t\mid\mathbf{x}_{t-1}),
  \qquad
  q(\mathbf{x}_t\mid\mathbf{x}_{t-1}) = \N\!\bigl(\mathbf{x}_t\mid\sqrt{1-\beta_t}\,\mathbf{x}_{t-1},\,\beta_t\mathbf{I}\bigr), \tag{6}
\]
and that marginalising over $\mathbf{x}_{1:t-1}$ gives
\[
  q(\mathbf{x}_t\mid\mathbf{x}_0) = \N\!\bigl(\mathbf{x}_t\mid\sqrt{\bar\alpha_t}\,\mathbf{x}_0,\,(1-\bar\alpha_t)\mathbf{I}\bigr),
  \qquad \bar\alpha_t = \prod_{\tau=1}^t (1-\beta_\tau). \tag{7}
\]

\textbf{Question C.1:} What is $q\!\bigl(\mathbf{x}_2\mid\mathbf{x}_0=[4,-4]^\top\bigr)$?

\begin{itemize}[noitemsep,label={}]
  \item[\textbf{A}] $\square$ \quad $\N\!\bigl(\mathbf{x}_2\mid [2,-2]^\top,\,\tfrac{3}{4}\mathbf{I}\bigr)$
  \item[\textbf{B}] $\square$ \quad $\N\!\bigl(\mathbf{x}_2\mid [2,-2]^\top,\,\tfrac{1}{4}\mathbf{I}\bigr)$
  \item[\textbf{C}] $\square$ \quad $\N\!\bigl(\mathbf{x}_2\mid [1,-1]^\top,\,\tfrac{3}{4}\mathbf{I}\bigr)$
  \item[\textbf{D}] $\square$ \quad $\N\!\bigl(\mathbf{x}_2\mid [1,-1]^\top,\,\tfrac{1}{4}\mathbf{I}\bigr)$
\end{itemize}

Give a brief argument for your answer.

\vspace{6cm}

\newpage
\section*{Part II \quad Representation learning}

\subsection*{Exercise D --- Pull-back metrics}

Consider the generative model $f:\mathbb{R}^d\to\mathbb{R}^D$ ($D\geq d$)
\[
  f(\mathbf{x}) = W_2\,\mathrm{ReLU}(W_1\mathbf{x}+\mathbf{b}_1)+\mathbf{b}_2, \tag{8}
\]
where $W_1\in\mathbb{R}^{D\times d}$ and $W_2\in\mathbb{R}^{D\times D}$ are full-rank matrices,
$\mathbf{b}_1\in\mathbb{R}^D$, $\mathbf{b}_2\in\mathbb{R}^D$, and
$\mathrm{ReLU}(x)=\max(0,x)$ is applied element-wise.  Define
\[
  \Lambda_1 = \mathrm{diag}\!\bigl(\mathbf{1}[W_1\mathbf{x}+\mathbf{b}_1 > \mathbf{0}]\bigr),\quad
  \Lambda_0 = \mathrm{diag}\!\bigl(\mathbf{1}[W_1\mathbf{x} > \mathbf{0}]\bigr), \tag{9}
\]
where $\mathbf{1}[\cdot]$ is the element-wise indicator function.

The pull-back metric is $G_{\mathbf{x}} = J_{\mathbf{x}}^\top J_{\mathbf{x}}$, where $J_{\mathbf{x}}$
is the Jacobian of $f$.  Which one of the following is correct?

\begin{itemize}[noitemsep,label={}]
  \item[\textbf{A}] $\square$ \quad $G_{\mathbf{x}} = W_1^\top \Lambda_1 W_2^\top W_2 \Lambda_1 W_1$
  \item[\textbf{B}] $\square$ \quad $G_{\mathbf{x}} = W_2^\top \Lambda_1 W_1^\top W_1 \Lambda_1 W_2$
  \item[\textbf{C}] $\square$ \quad $G_{\mathbf{x}} = W_1^\top \Lambda_0 W_2^\top W_2 \Lambda_0 W_1$
  \item[\textbf{D}] $\square$ \quad $G_{\mathbf{x}} = W_2^\top \Lambda_0 W_1^\top W_1 \Lambda_0 W_2$
\end{itemize}

Give a brief argument for your answer.

\vspace{5cm}

\newpage
\subsection*{Exercise E --- Curve energy}

Consider the generative model $\mathbf{y}=f(\mathbf{x})$, where
\[
  f(\mathbf{x}) = \bigl[x_1^2,\; x_1 x_2,\; x_2\bigr]^\top. \tag{10}
\]
Furthermore, consider the two latent curves
\[
  \mathbf{c}_1(t) = t[1,0]^\top \qquad\text{and}\qquad \mathbf{c}_2(t) = t[0,1]^\top, \quad t\in[0,1].
\]

\textbf{Question E.1:} Derive the pull-back metric
$G_{\mathbf{x}} = J_{\mathbf{x}}^\top J_{\mathbf{x}}$, where $J_{\mathbf{x}}$ is the Jacobian of
$f$ at $\mathbf{x}$.

\vspace{5cm}

\textbf{Question E.2:} The energy of a curve is defined as
\[
  E(\mathbf{c}) = \int_0^1 \dot{\mathbf{c}}(t)^\top G_{\mathbf{c}(t)}\,\dot{\mathbf{c}}(t)\,\mathrm{d}t, \tag{11}
\]
where $\dot{\mathbf{c}}(t)=\mathrm{d}\mathbf{c}(t)/\mathrm{d}t$.  Evaluate $E(\mathbf{c}_1)$ and
$E(\mathbf{c}_2)$.

\[
  E(\mathbf{c}_1) = \hspace{6cm} E(\mathbf{c}_2) =
\]

Sketch your derivations below.

\vspace{5cm}

\newpage
\subsection*{Exercise F --- Riemannian integration}

Consider the Riemannian metric over $\mathbf{x}=(x_1,x_2)\in[0,1]\times[0,1]$
\[
  G_{\mathbf{x}} = \mathrm{diag}(x_2^2,\; x_1^2). \tag{12}
\]

\textbf{Question F.1:} Derive the volume measure $\mathrm{d}M(\mathbf{x})=\sqrt{\det G_{\mathbf{x}}}\,\mathrm{d}\mathbf{x}$.

\[
  \mathrm{d}M(\mathbf{x}) =
\]

Sketch your derivations below.

\vspace{3cm}

\textbf{Question F.2:} Let $\Omega=[0,1]\times[0,1]$.  Evaluate the integrals
\[
  I_1 = \int_\Omega x_1^2\,\mathrm{d}M(\mathbf{x}), \qquad
  I_2 = \int_\Omega x_2\,\mathrm{d}M(\mathbf{x}).
\]

\[
  I_1 = \hspace{6cm} I_2 =
\]

Sketch your derivations below.

\vspace{5cm}

\newpage
\subsection*{Exercise G --- ELBO and evidence}

Recall that for a latent variable model with prior $p(\mathbf{z})$, likelihood $p(\mathbf{x}\mid\mathbf{z})$,
and approximate posterior $q(\mathbf{z}\mid\mathbf{x})$, the ELBO is defined as
\[
  \ELBO(\mathbf{x}) = \E_{q(\mathbf{z}\mid\mathbf{x})}\!\Bigl[\ln\frac{p(\mathbf{x},\mathbf{z})}{q(\mathbf{z}\mid\mathbf{x})}\Bigr]. \tag{13}
\]

\textbf{Question G.1:} Show that the ELBO can be written as
\[
  \ELBO(\mathbf{x}) = \ln p(\mathbf{x}) - \KL\!\bigl[q(\mathbf{z}\mid\mathbf{x})\;\|\;p(\mathbf{z}\mid\mathbf{x})\bigr]. \tag{14}
\]

\emph{Hint:} Use the factorisation $p(\mathbf{x},\mathbf{z}) = p(\mathbf{z}\mid\mathbf{x})\,p(\mathbf{x})$.

\vspace{7cm}

\newpage
\section*{Part III \quad Graph models}

\subsection*{Exercise H --- Graph Laplacians}

The random-walk graph Laplacian is defined as $L_{\mathrm{rw}} = I - D^{-1}A$, where $A$ is the
adjacency matrix and $D$ the degree matrix.  Consider a graph with
\[
  L_{\mathrm{rw}} =
  \begin{bmatrix}
    1  & -1   & 0    & 0    & 0    \\
    -\tfrac12 & 1 & -\tfrac12 & 0 & 0 \\
    0  & -\tfrac12 & 1 & -\tfrac12 & 0 \\
    0  &  0  & -\tfrac12 & 1 & -\tfrac12 \\
    0  &  0  & 0    & -1   & 1
  \end{bmatrix},
\]
with nodes labelled $a$--$e$ in row order.

\textbf{Question H.1:} Sketch the graph.

\vspace{3cm}

\textbf{Question H.2:} Consider a random walk on this graph starting at node $a$.  What is the
probability $p$ of being at node $b$ after exactly 3 steps?

\[
  p =
\]

Give a brief argument for your answer.

\vspace{3cm}

\subsection*{Exercise I --- Weisfeiler-Lehman test}

Consider the following graph with nodes $a, b, c, d$: a triangle formed by $a$--$b$--$c$--$a$,
with an additional edge $a$--$d$ (so $d$ is a pendant connected only to $a$).  Run the
Weisfeiler-Lehman (WL) algorithm:
\begin{enumerate}[noitemsep]
  \item Assign initial labels $\ell^{(0)}_v = d_v$ (degree of node $v$).
  \item Iteratively: $\ell^{(i)}_v = \mathrm{hash}\!\bigl(\{\!\{\ell^{(i-1)}_u : u\in\mathcal{N}(v)\}\!\}\bigr)$.
\end{enumerate}

\textbf{Question I.1:} Which of the following are true?  Select all that apply.

\begin{itemize}[noitemsep,label={}]
  \item[\textbf{A}] $\square$ \quad $\ell^{(0)}_a = 3$
  \item[\textbf{B}] $\square$ \quad $\ell^{(0)}_d = 2$
  \item[\textbf{C}] $\square$ \quad $\ell^{(1)}_b = \ell^{(1)}_c$
  \item[\textbf{D}] $\square$ \quad $\ell^{(1)}_a = \ell^{(1)}_b$
  \item[\textbf{E}] $\square$ \quad $\ell^{(1)}_d = \ell^{(1)}_b$
\end{itemize}

\textbf{Question I.2:} After how many rounds does the WL algorithm halt for this graph?

\[
  \text{Number of rounds} =
\]

Give a brief argument.

\vspace{3cm}

\newpage
\subsection*{Exercise J --- Shallow embeddings}

We use the squared-distance decoder $\mathrm{DEC}(\mathbf{z}_u,\mathbf{z}_v) = \|\mathbf{z}_u-\mathbf{z}_v\|^2$
in a 2-dimensional embedding space with $\mathbf{z}_v=[x,y]^\top$.

\textbf{Question J.1:} Let $\mathbf{z}_u=[0,1]^\top$.  Where could node $v$ be embedded if the
decoder evaluates to $\mathrm{DEC}(\mathbf{z}_u,\mathbf{z}_v)=0.25$?

\begin{itemize}[noitemsep,label={}]
  \item[\textbf{A}] $\square$ \quad On a circle of radius $0.5$ centred at $\mathbf{z}_u$,
    i.e.\ $x^2+(y-1)^2=0.25$
  \item[\textbf{B}] $\square$ \quad On a circle of radius $1$ centred at the origin,
    i.e.\ $x^2+y^2=1$
  \item[\textbf{C}] $\square$ \quad On the line $x=0$
  \item[\textbf{D}] $\square$ \quad On the line $y=0$
  \item[\textbf{E}] $\square$ \quad None of the above
\end{itemize}

Show the main steps of your derivation.

\vspace{4cm}

\subsection*{Exercise K --- Message passing}

Consider a star graph with centre $a$ and leaves $b,c,d$ (so $\mathcal{N}(a)=\{b,c,d\}$ and
$\mathcal{N}(b)=\mathcal{N}(c)=\mathcal{N}(d)=\{a\}$).  We use the message-passing GNN
\[
  h_u^{(k)} = w_0 + w_1 h_u^{(k-1)} + w_2 \sum_{v\in\mathcal{N}(u)} h_v^{(k-1)}, \tag{15}
\]
with $w_0=w_1=w_2=1$ and initial features $h_u^{(0)}=1$ for all $u$.

\textbf{Question K.1:} What are the node features after two message-passing rounds?

\[
  h_a^{(2)} = \qquad\qquad h_b^{(2)} = h_c^{(2)} = h_d^{(2)} =
\]

Show the main steps of your derivation.

\vspace{3cm}

\textbf{Question K.2:} A single GNN layer can be written as the matrix--vector product
$\mathbf{h}^{(k)} = \mathbf{w}_0 + M\,\mathbf{h}^{(k-1)}$, where
$\mathbf{h}^{(k)}=[h_a^{(k)},h_b^{(k)},h_c^{(k)},h_d^{(k)}]^\top$ and
$\mathbf{w}_0=[w_0,w_0,w_0,w_0]^\top$.
Again with $w_1=w_2=1$, fill in all non-zero entries of $M$.

\[
  M = \begin{bmatrix} \phantom{1} & \phantom{1} & \phantom{1} & \phantom{1} \\[6pt]
                      & & & \\[6pt]
                      & & & \\[6pt]
                      & & & \end{bmatrix}
\]

\end{document}
